

% Anpassung an Landessprache
\usepackage[ngerman]{babel}
 
% Verwenden von Sonderzeichen und Silbentrennung
\usepackage[utf8]{inputenc}	
\usepackage[T1]{fontenc}			
\usepackage{textcomp} 																% Euro-Zeichen und andere
\usepackage[babel,german=quotes]{csquotes}						% Anführungszeichen
\RequirePackage[ngerman=ngerman-x-latest]{hyphsubst} 	% erweiterte Silbentrennung

% Seitenzahlen nach Kudraß: oben am Rand, nur die Zahl, keine Zahl auf der
% ersten Seite eines Kapitels bzw. Verzeichnisses und auf dem Abstract
\usepackage{scrlayer-scrpage}
\clearpairofpagestyles
\ohead{\pagemark}
\pagestyle{scrheadings}
\renewcommand*{\chapterpagestyle}{empty}

% Befehle aus AMSTeX für mathematische Symbole z.B. \boldsymbol \mathbb
\usepackage{amsmath,amsfonts}

% Schrift Times (Kudraß: 12 pt Times oder 11 pt Arial); newtxmath nach amsmath.
% newtxtext lädt xcolor, deshalb die Farbnamen vorher als Option übergeben
\PassOptionsToPackage{dvipsnames}{xcolor}
\usepackage{newtxtext,newtxmath}

% Zeilenabstände und Seitenränder 
\usepackage{setspace}
\usepackage{geometry}

% Einbinden von JPG-Grafiken
\usepackage{graphicx}
\graphicspath{{Abbildungen/}}

% zum Umfließen von Bildern
% Verwendung unter http://de.wikibooks.org/wiki/LaTeX-Kompendium:_Baukastensystem#textumflossene_Bilder
\usepackage{floatflt}

% Verwendung von vordefinierten Farbnamen zur Colorierung
% Palette und Verwendung unter http://kitt.cl.uzh.ch/kitt/CLinZ.CH/src/Kurse/archiv/LaTeX-Kurs-Farben.pdf
% (xcolor statt color, weil newtxtext und TikZ ohnehin xcolor laden)
\usepackage[dvipsnames]{xcolor}

% eigene Diagramme direkt in LaTeX (pstricks funktioniert mit pdfLaTeX nicht)
\usepackage{tikz}
\usetikzlibrary{positioning,arrows.meta,shapes,fit,calc,shadows}

% Tabellen
\usepackage{array}
\usepackage{longtable}

% einfache Grafiken im Code
% Einführung unter http://www.math.uni-rostock.de/~dittmer/bsp/pstricks-bsp.pdf
\usepackage{pstricks}

% Quellcodeansichten
\usepackage{verbatim}
\usepackage{moreverb} 											% für erweiterte Optionen der verbatim Umgebung
% Befehle und Beispiele unter http://www.ctex.org/documents/packages/verbatim/moreverb.pdf
\usepackage{listings} 											% für angepasste Quellcodeansichten siehe
% Kurzeinführung unter http://blog.robert-kummer.de/2006/04/latex-quellcode-listing.html

% Glossar und Abbildungsverzeichnis
\usepackage[
nonumberlist, %keine Seitenzahlen anzeigen
acronym,      %ein Abkürzungsverzeichnis erstellen
toc          %Einträge im Inhaltsverzeichnis
]      %im Inhaltsverzeichnis auf section-Ebene erscheinen
{glossaries}

% verlinktes und Farblich angepasstes Inhaltsverzeichnis
\usepackage[pdftex,
colorlinks=true,
linkcolor=InterneLinkfarbe,
citecolor=InterneLinkfarbe, % Zitate wie interne Verweise statt grün
urlcolor=ExterneLinkfarbe]{hyperref}
\usepackage[all]{hypcap}

% URL verlinken, lange URLs umbrechen
\usepackage{url}

% sorgt dafür, dass Leerzeichen hinter parameterlosen Makros nicht als Makroendezeichen interpretiert werden
\usepackage{xspace}

% Beschriftungen für Abbildungen und Tabellen
\usepackage{caption}

% Entwicklerwarnmeldungen entfernen
\usepackage{scrhack}